\documentclass[12pt]{article}
\usepackage{graphicx,amsmath,amsfonts,amssymb,epsfig,euscript,enumerate}
\usepackage[T1]{fontenc}
\usepackage[utf8x]{inputenc}

\newtheorem{exercise}{Ejercicio}
\newcommand{\bej}{\begin{exercise}\rm}
\newcommand{\fej}{\end{exercise}}

\newcommand{\R}{\mathbb{R}}
\newcommand{\C}{\mathbb{C}}
\def\dt{\Delta t}
\def\dx{\Delta x}

\topmargin-2cm \vsize 29.5cm \hsize 21cm
\setlength{\textwidth}{16.75cm}\setlength{\textheight}{23.5cm}
\setlength{\oddsidemargin}{0.0cm}
\setlength{\evensidemargin}{0.0cm}

\begin{document}
\centerline{{\small Universidad de Buenos Aires - Facultad de Ciencias Exactas y Naturales - Depto. de Matemática}}
 
 \vskip 0.2cm
 \hrulefill
 \vskip 0.2cm

 \centerline{{\bf\Huge {\sc Elementos de Cálculo Numérico}}}
 \vskip 0.2cm
 \centerline{\ttfamily Primer Cuatrimestre 2026}
 \hrulefill

 \bigskip
 \centerline{\bf Ejercicios tipo parcial -- Guías 1 a 3}
 \bigskip

\begin{exercise}[Ajuste racional por cuadrados m\'inimos y SVD]
Dados datos $(x_i,f_i)$, se propone ajustar el modelo racional
\[
r(x)=\frac{a+bx}{1+cx}
\]
por cuadrados m\'inimos.

\begin{enumerate}[(a)]
\item Muestre que, partiendo de la relaci\'on aproximada $f_i\approx r(x_i)$ y multiplicando por $(1+cx_i)$, el problema puede escribirse como un problema lineal de cuadrados m\'inimos para los par\'ametros $a,b,c$.

\item Escriba el sistema correspondiente en forma matricial $A\theta\approx b$ con $\theta=(a,b,c)^T$, indicando expl\'icitamente las filas de $A$ y el vector de datos del segundo miembro.

\item Proponga un algoritmo num\'ericamente estable para resolverlo usando la SVD de $A$, y describa c\'omo recuperar luego la funci\'on ajustada $r(x)$.
\end{enumerate}
\end{exercise}

\begin{exercise}[Integral con singularidad algebraica integrable]
Considere
\[
I=\int_0^1 t^{\alpha}f(t)\,dt, \qquad -1<\alpha<0,
\]
con $f$ continua en $[0,1]$.

\begin{enumerate}[(a)]
\item Discuta si es posible aplicar directamente la regla del trapecio compuesta al integrando $t^{\alpha}f(t)$ y justifique su respuesta.

\item Muestre que sumando y restando el valor $f(0)\int_0^1 t^{\alpha}\,dt$, la integral puede descomponerse en un término regularizado y un término singular que se puede calcular explícitamente. A partir de esta descomposición, plantee un esquema numérico basado en la regla de trapecios compuesta para aproximar $I$ y muestre que el orden de convergencia del método es al menos $O(h)$.
\end{enumerate}
\end{exercise}

\begin{exercise}[Transformada de Haar]
En este ejercicio se introduce la transformada de Haar para vectores de longitud $n=2^m$.

Dado $x\in\mathbb{R}^n$, en un nivel se forman promedios y diferencias por pares:
\[
u_j=\frac{x_{2j}+x_{2j+1}}{\sqrt2},
\qquad
v_j=\frac{x_{2j}-x_{2j+1}}{\sqrt2},
\qquad j=0,\ldots,\frac n2-1.
\]
La salida de ese nivel es el vector
\[
\bigl(u_0,\ldots,u_{n/2-1},v_0,\ldots,v_{n/2-1}\bigr)^T,
\]
y luego el mismo procedimiento se aplica recursivamente al bloque de promedios $(u_j)$ hasta llegar a longitud $1$.

\begin{enumerate}[(a)]
\item Escriba pseudocódigo para aplicar la transformada de Haar a un vector $x\in\mathbb{R}^n$ usando la definición recursiva anterior y calcule la complejidad computacional.

\item Utilizando que los pasos de promedios y diferencias por pares pueden escribirse como la multiplicación por las matrices $A_n$ y $B_n$ respectivamente, dadas por
\[
A_n=\frac1{\sqrt2}\begin{pmatrix}I_{n/2} & I_{n/2}\end{pmatrix},
\qquad
B_n=\frac1{\sqrt2}\begin{pmatrix}I_{n/2} & -I_{n/2}\end{pmatrix},
\]
y que la transformada de Haar total se puede escribir como la multiplicación por la matriz $H_n$ dada por
\[    
H_n=\begin{pmatrix}
H_{n/2}A_n\\
B_n
\end{pmatrix}.
\]
demuestre por inducción que $H_n^TH_n=I_n$, es decir, que la transformada de Haar es una transformación ortogonal.
\item Utilizando el resultado anterior, ¿Cuál es el condicionamiento de la matriz $H_n$? ¿Qué nos dice esto sobre la estabilidad numérica de la transformada de Haar? Justifique su respuesta.
\end{enumerate}
\end{exercise}


\begin{exercise}[Función de Green cuasi-periódica y FFT] En problemas de difracción por estructuras periódicas, la función de Green cuasi-periódica se puede representar mediante una serie espectral. Para fines computacionales, a menudo se utiliza una versión truncada de esta serie:
\[
G^{\mathrm{qp}}_M(x,y)=\frac{i}{2d}\sum_{m=-M}^{M}
\frac{\exp\!\bigl(i\alpha_m x+i\beta_m|y|\bigr)}{\beta_m},
\qquad
\alpha_m=\alpha+\frac{2\pi m}{d},
\]
donde $\beta_m=\sqrt{k^2-\alpha_m^2}$ (rama con $\operatorname{Im}(\beta_m)\ge 0$), $k>0$ y $\alpha\in\mathbb{R}$ son datos.

Sea una grilla uniforme en $x$ dada por $x_j=j\,d/N$, $j=0,\ldots,N-1$, y suponga $N\ge 2M+1$.

\begin{enumerate}[(a)]
\item Reescriba $G^{\mathrm{qp}}_M(x_j,y)$ en forma de suma trigonométrica
\[
G^{\mathrm{qp}}_M(x_j,y)=e^{i\alpha x_j}\sum_{m=-M}^{M} c_m(y)\,e^{2\pi i mj/N},
\]
identificando explícitamente los coeficientes $c_m(y)$.

\item A partir de la expresión anterior, diseñe un algoritmo basado en FFT para calcular simultáneamente los valores $\{G^{\mathrm{qp}}_M(x_j,y)\}_{j=0}^{N-1}$ para un $y$ fijo.

\item Indique cómo extender el procedimiento para varios valores $y_1,\ldots,y_q$. ¿Qué magnitudes puede pre-calcular y reutilizar para obtener un algoritmo lo más eficiente posible?¿Cuál sería la complejidad computacional total del algoritmo para calcular $G^{\mathrm{qp}}_M(x_j,y_k)$ para $j=0,\ldots,N-1$ y $k=1,\ldots,q$? Justifique su respuesta.
\end{enumerate}
\end{exercise}


\end{document}
